\documentclass[12pt]{article}

% Packages
\usepackage{amsfonts}
\usepackage{amsmath}
\usepackage{amssymb}
\usepackage{amsthm}
\usepackage{bm}
\usepackage{bbm}
\usepackage[style=numeric]{biblatex}
\usepackage{caption}
\usepackage{cancel}
\usepackage{empheq}
\usepackage{fancyhdr}
\usepackage[margin=1in]{geometry}
\usepackage{hyperref}
\usepackage[nameinlink]{cleveref}
\usepackage{mathtools}
\usepackage{pgfplots}
\usepackage{physics}
\usepackage{siunitx}
\usepackage{tensor}
\usepackage{tikz}
\usepackage{upgreek}
\usepackage{wasysym}

\usetikzlibrary{calc,patterns,angles,quotes}
\usetikzlibrary{decorations.markings}

% PGFPlots
\pgfplotsset{compat=newest} 

% Variables
\newcommand{\theAuthor}{Paul Virally}
\newcommand{\theTitle}{Multivariate Gaussians}
\newcommand{\theSubtitle}{}
\newcommand{\theDate}{\today}

% Headers and footers
\pagestyle{fancy}
\lhead{\theAuthor{}}
\chead{\textit{\theSubtitle{}}}
\rhead{\theTitle{}}
\lfoot{}
\cfoot{}
\rfoot{Page \thepage}

\newcommand\red[1]{
    \ensuremath{\textcolor{red}{#1}}
}

\newcommand\blue[1]{
    \ensuremath{\textcolor{blue}{#1}}
}

\newcommand\todo[1]{
    \Huge{\color{red}{TODO:~#1}}\normalsize{}
}

\newcommand{\kbar}{\mathchar'26\mkern-9mu k}

\renewcommand\qedsymbol{\(\blacksquare\)}
\newtheorem{theorem}{Theorem}
\numberwithin{theorem}{section}
\numberwithin{equation}{section}
\numberwithin{figure}{section}

\newcommand*\circled[1]{\tikz[baseline=(char.base)]{\node[shape=circle,draw,inner sep=2pt] (char) {#1};}}

% \addbibresource{bib.bib}

% Document
\begin{document}

% Cover page
% \begin{titlepage}
%   \begin{center}
%       \textsc{\Huge{\theTitle}}\\[0.5cm]
%       \textsc{\Large{\theSubtitle}}\\[4cm]
%       \theAuthor{}\\[0.2cm]
%       \theDate{}
%   \end{center}
% \newpage
% \end{titlepage}
% \newpage

\tableofcontents
\newpage

\section{Easy Multivariate Gaussians}

The general multivariate Gaussian probability density function, \(f : \mathbb{R}^{d} \to \mathbb{R}\) is given by:
\begin{align*}
    f\qty(\ket{x}) &= \frac{1}{\sqrt{\qty(2\pi)^{d}\abs{\bm{\Sigma}}}} \exp(-\frac{1}{2}\qty(\bra{x} - \bra{\mu})\bm{\Sigma}^{-1}\qty(\ket{x} - \ket{\mu}))\stepcounter{equation}\tag{\theequation}\label{eq:pdf}
\end{align*}
where \(\ket{\mu}\in \mathbb{R}^{d}\) is the mean vector and \(\bm{\Sigma}\succ \bm{0}\) is the covariance matrix. In our case, \(\bm{\Sigma}\) takes on a very particular form:
\begin{align*}
    \bm{\Sigma}\qty(r) &= \qty(1 - r)\bm{\mathbbm{1}} + r \op{1} \\
                       &= \begin{bmatrix}
                            1 & r & \cdots & r \\
                            r & 1 & \cdots & r \\
                            \vdots & \vdots & \ddots & \vdots \\
                            r & r & \cdots & 1
                        \end{bmatrix}
\end{align*}
where \(\ket{1}\) is a vector of all ones. This simple structure gives us a lot. First, notice that \(\ket{1}\) is an eigenvector of \(\bm{\Sigma}\):
\begin{align*}
    \bm{\Sigma}\ket{1} &= \qty(1 - r)\bm{\mathbbm{1}}\ket{1} + r\ket{1}\underbrace{\bra{1}\ket{1}}_{d} \\
                       &= (1 + r\qty(d-1))\ket{1}
\end{align*}
Now consider any vector \(\ket{x}\in \mathbb{R}^{d}\) orthogonal to \(\ket{1}\),
i.e. such that \(\bra{1}\ket{x} = 0\). Then:
\begin{align*}
    \bm{\Sigma}\ket{x} &= \qty(1 - r)\bm{\mathbbm{1}}\ket{x} + \ket{1}\underbrace{\bra{1}\ket{x}}_{0} \\
                       &= (1 - r)\ket{x}
\end{align*}
The span of vectors orthogonal to \(\ket{1}\) is of dimension \(d - 1\). Thus, the eigenvalues of \(\bm{\Sigma}\) are \(1+r\qty(d-1)\) with multiplicity 1 and \(1 - r\) with multiplicity \(d - 1\). The determinant of \(\bm{\Sigma}\) is then:
\begin{align*}
    \abs{\bm{\Sigma}} &= (1 + r\qty(d-1))(1 - r)^{d - 1} \stepcounter{equation}\tag{\theequation}\label{eq:det}
\end{align*}
For \(\bm{\Sigma}\) to be positive definite, we require that both eigenvalues are positive, which gives us the range:
\begin{align*}
    -\frac{1}{d - 1} < r < 1\stepcounter{equation}\tag{\theequation}\label{eq:range}
\end{align*}
From the Sherman-Morrison formula, we can also compute the inverse of \(\bm{\Sigma}\):
\begin{align*}
    \bm{\Sigma}^{-1} &= \qty(1-r)^{-1}\bm{\mathbbm{1}} - \frac{\qty(1-r)^{-1}r\op{1}\qty(1-r)^{-1}}{1 + \bra{1}\qty(1-r)^{-1}r\ket{1}} \\
                     &= \qty(1-r)^{-1}\bm{\mathbbm{1}} - \frac{r\qty(1-r)^{-2}}{1 + rd\qty(1-r)^{-1}}\op{1}\\
                     &= \qty(1-r)^{-1}\qty[\bm{\mathbbm{1}} - \frac{r}{1 + r\qty(d - 1)}\op{1}]\stepcounter{equation}\tag{\theequation}\label{eq:inv}
\end{align*}
Let \(\ket{\tilde{x}} = \ket{x} - \ket{\mu}\), then the quadratic form in the exponent of~\cref{eq:pdf} can be written as:
\begin{align*}
    \bra{\tilde{x}}\bm{\Sigma}^{-1}\ket{\tilde{x}} &= \qty(1-r)^{-1}\qty[\bra{\tilde{x}}\ket{\tilde{x}} - \frac{r}{1 + r\qty(d - 1)}\bra{\tilde{x}}\ket{1}\bra{1}\ket{\tilde{x}}] \\
                                                           &= \qty(1-r)^{-1}\qty[\bra{\tilde{x}}\ket{\tilde{x}} - \frac{r}{1 + r\qty(d - 1)}\qty(\sum_{i=1}^{d}\tilde{x}_i)^2]\stepcounter{equation}\tag{\theequation}\label{eq:quad}
\end{align*}
The multivariate Gaussian distribution is then:
\begin{align*}
    f\qty(\ket{x}) &= \frac{1}{\sqrt{2\pi(1 + r\qty(d-1))(1 - r)^{d - 1}}} \exp\qty(-\frac{1}{2(1-r)}\qty[\norm{\ket{x} - \ket{\mu}}_2^2 - \frac{r}{1 + r\qty(d - 1)}\qty(\sum_{i=1}^{d}\qty(x_i - \mu_i))^2])\stepcounter{equation}\tag{\theequation}\label{eq:pdf-simple}
\end{align*}

\section{Tensors and the Tucker HOSVD}

Let \(H_i\) be a Hilbert space of dimension \(d_i\). An \(n\)-way tensor
\(\mathbb{T}\) is an element of the tensor product \(H = H_1 \otimes H_2 \otimes
\cdots \otimes H_N\) with coordinates \(T_{i_1, i_2,\ldots, i_N}\) in some bases
of the \(H_i\). The tensor itself acts on elements of the dual spaces
\(H_i^{*}\) (linear functionals, or bras) as follows:
\begin{align*}
    \mathbb{T}\qty(\bra*{v_1}, \bra*{v_2}, \ldots, \bra*{v_N}) &= \sum_{i_1, \ldots, i_N} T_{i_1, i_2,\ldots, i_N} \bra*{v_1}\ket*{i_1} \bra*{v_2}\ket*{i_2} \cdots \bra*{v_N}\ket*{i_N}
\end{align*}
where \(\ket*{i_k}\) is the \(i_k\)-th basis vector of \(H_k\).

For tensors, apparently there is no unique analogue of the SVD for matrices. In
an SVD, you get the optimal rank-\(r\) approximation of a matrix by keeping the
\(r\) largest singular values while simultaneously getting orthonormal bases for
the row and column spaces. For tensors, you can get the optimal rank-\(r\)
approximation by keeping the \(r\) largest singular values, but you don't get
orthonormal bases for the spaces on which the tensor acts. You can also get
orthonormal bases, but then you don't get the optimal rank-\(r\) approximation.
The Tucker Higher-order singular value decomposition (HOSVD) is the latter: a
decomposition of a tensor that focuses on getting orthonormal bases for the
spaces on which the tensor acts, but there is no guarantee that you get the
optimal rank-\(r\) approximation.

The HOSVD relies on ``matricization'' (mode unfolding, or \texttt{reshape}) of
the tensor. Define \(\mathbb{T}_{\qty(k)} : H_k \to \bigotimes_{j\neq k} H_j\)
as the mode-\(k\) unfolding of \(\mathbb{T}\). In some basis, this is just a
\texttt{reshape} of the coordinates of \(\mathbb{T}\) (\(T_{i_1, \ldots, i_N}\))
into a matrix \(T_{(k)} \in \mathbb{C}^{d_k \times \prod_{j\neq k} d_j}\).

To form the HOSVD, one computes the SVD of each mode unfolding \(T_{(k)}\):
\begin{align*}
    T_{(k)} &= U_k \Sigma_k V_k^\dagger
\end{align*}
where, as usual, \(U_k\) and \(V_k\) are unitary and \(\Sigma_k\) is diagonal
with non-negative entries. From here, the core tensor \(\mathbb{S}\) is defined
as:
\begin{align*}
    \mathbb{S} &= \mathbb{T} \times_1 V_1^\dagger \times_2 V_2^\dagger \times_3 \cdots \times_N V_N^\dagger
\end{align*}
where \(\times_k\) is the mode-\(k\) product of a tensor by a matrix. The HOSVD
of \(\mathbb{T}\) is then given by:
\begin{align*}
    \mathbb{T} &= \mathbb{S} \times_1 V_1 \times_2 V_2 \times_3 \cdots \times_N V_N
\end{align*}

\section{HOSVD of Multivariate Gaussians}
We found the multivariate Gaussian distribution \(f\qty(\ket{x})\)
in~\cref{eq:pdf-simple}. We can think of this function as the coordinates of a
tensor \(\mathbb{F}\):
\begin{align*}
    \mathbb{F}\qty(\bra{g}) &= \int_{\mathbb{R}^d} f\qty(\ket{x}) g_1\qty(x_1) \cdots g_n\qty(x_n) \dd[n]{\ket{x}}
\end{align*}
Note that the Hilbert spaces in question here are \(H_i = L^2\qty(\mathbb{R})\)
for \(i = 1, \ldots, n\).

Let's compute the mode-\(k\) unfolding of this tensor, \(F_{\qty(k)} : H_k \to
\bigotimes_{j \neq k}H_j\). Let \(\ket{x} = \qty(x_k, \ket{x_{\kbar}})\) where
\(\ket{x_{\kbar}}\) is the vector of all coordinates except \(x_k\). In a basis,
this looks like:
\begin{align*}
    \qty(F_{\qty(k)}g)\qty(\ket{x_{\kbar}}) &= \int_{\mathbb{R}} f\qty(x_k, \ket{x_{\kbar}}) g\qty(x_k) \dd{x_k}
\end{align*}
In the HOSVD, all we really need are the right singular vectors and singular
values of \(F_{\qty(k)}\). The right singular vectors of \(F_{\qty(k)}\) are the
eigenvectors of \(F_{\qty(k)}^\dagger F_{\qty(k)}\) and the singular values are
the square roots of the corresponding eigenvalues. Let's compute
\(F_{\qty(k)}^\dagger F_{\qty(k)}\):
\begin{align*}
    \qty(C_kg)\qty(t) &= \qty(F_{\qty(k)}^\dagger\qty(F_{k}g))\qty(t) \\
                      &= \int_{\mathbb{R}^{d-1}} f\qty(t, \ket{x_{\kbar}}) \qty(F_{\qty(k)}g)\qty(\ket{x_{\kbar}}) \dd[n-1]{\ket{x_{\kbar}}} \\
                      &= \int_{\mathbb{R}^{d-1}} f\qty(t, \ket{x_{\kbar}}) \qty(\int_{\mathbb{R}} f\qty(s, \ket{x_{\kbar}})g\qty(s)\dd{s}) \dd{x_k}) \dd[n-1]{\ket{x_{\kbar}}} \\
                      &= \int_{\mathbb{R}} \underbrace{\qty(\int_{\mathbb{R}^{d-1}} f\qty(t, \ket{x_{\kbar}})f\qty(s, \ket{x_{\kbar}}))}_{\kappa_{\kbar}\qty(t, s)} g\qty(s) \dd{s}
\end{align*}
where \(\kappa_{\kbar}\qty(t, s)\) is the kernel of the integral operator
\(C_k\). The right singular vectors of \(F_{\qty(k)}\) are the functions
\(v_{k}^{\qty(n)}\) such that \(\displaystyle\qty(C_kv_k^{\qty(n)})\qty(s) =
\int_{\mathbb{R}} \kappa_{\kbar}\qty(t, s) v_k^{\qty(n)}\qty(t) \dd{t} =
\lambda_k^{\qty(n)}v_k^{\qty(n)}\) and the singular values are
\(\sigma_k^{\qty(n)} = \sqrt{\lambda_k^{\qty(n)}}\).

Let's compute the kernel \(\kappa_{\kbar}\qty(t, s)\):
\begin{align*}
    \kappa_{\kbar}\qty(t, s) &= \int_{\mathbb{R}^{d-1}} f\qty(t, \ket{x_{\kbar}})f\qty(s, \ket{x_{\kbar}}) \dd[n-1]{\ket{x_{\kbar}}}
\end{align*}
First, let's define some variables to make the expression for \(f\) more
compact:
\begin{align*}
    \text{Let }\mathcal{N} &= \qty(\qty(2\pi)^{d}\qty(1 + r\qty(d - 1))\qty(1-r)^{d-1})^{-1 / 2} \\
    \text{Let } a &= \frac{1}{1-r} \\
    \text{Let } b &= \frac{r}{1 + r\qty(d-1)} \\
    \text{Let } \ket{\tilde{x}_{\kbar}} &= \ket{x_{\kbar}} - \ket{\mu} \\
    \text{Let } \tilde{\varsigma} &= \sum_{i=1}^{d}\ket{\tilde{x}_{\kbar}}_i \\
    \text{Let } \tilde{n} &= \norm{\ket{\tilde{x}_{\kbar}}}_2 \\
    \implies f(u, \ket{x_{\kbar}}) &= \mathcal{N} \exp\qty(-\frac{a}{2}\qty[\tilde{n}^2 + (u - \mu_k)^2 - b\qty(\tilde{\varsigma} + \qty(u - \mu_k))^2]) \\
    \implies f\qty(t, \ket{x_{\kbar}})f\qty(s, \ket{x_{\kbar}}) &= \mathcal{N}^2 \exp\bigg(\!\!-\frac{a}{2}\bigg[\tilde{n}^2 + (t - \mu_k)^2 - b\qty(\tilde{\varsigma} + \qty(t - \mu_k))^2 + \\
                             &\qquad + \tilde{n}^2 + (s - \mu_k)^2 - b\qty(\tilde{\varsigma} + \qty(s - \mu_k))^2\,\bigg]\bigg) \\
                             &= \mathcal{N}^2 \exp\qty(-\frac{a\qty(1 - b)}{2}\qty[\qty(t - \mu_k)^2 + \qty(s - \mu_k)^2]) \times \\
                             &\quad\times \exp(-a\qty[\tilde{n}^2 - b \tilde{\varsigma}^2 - 2\tilde{\varsigma}\qty(t + s - 2\mu_k)])
\end{align*}
The kernel is then:
\begin{align*}
    \kappa_{\kbar}\qty(t, s) &= \int_{\mathbb{R}^{d-1}} f\qty(t, \ket{x_{\kbar}})f\qty(s, \ket{x_{\kbar}}) \dd[n-1]{\ket{x_{\kbar}}} \\
                             &= \frac{\exp(-\frac{1 + r\qty(d - 2)}{2\qty(1-r)\qty[1 + r\qty(d-1)]}\qty[\qty(t - \mu_k)^2 + \qty(s - \mu_k)^2])}{\qty(2\pi)^{d}\qty(1 + r\qty(d-1))\qty(1-r)^{d-1}} \times \\
                             &\quad \times \underbrace{\int_{\mathbb{R}^{d-1}} \exp\qty(-a\qty[\tilde{n}^2 - b \tilde{\varsigma}^2 - 2\tilde{\varsigma}\qty(t + s - 2\mu_k)]) \dd[n-1]{\ket{x_{\kbar}}}}_{I}
\end{align*}
Focusing on \(I\), we can split the vector \(\ket{x_{\kbar}}\) into its
projection onto \(\ket{1}\) and its projection onto the subspace orthogonal to
\(\ket{1}\): \(\displaystyle \ket{x_{\kbar}} = \frac{\alpha}{\sqrt{\bra{1}\ket{1}}}
\ket{1} + \ket{\perp}\). This simplifies a
few of the quantities in \(I\):
\begin{align*}
    \tilde{n}^2 &= \bra{x_{\kbar}}\ket{x_{\kbar}} \\
                &= \frac{\alpha^2}{\bra{1}\ket{1}} \bra{1}\ket{1} + \bra{\perp}\ket{\perp} \\
                &= \alpha^2 + \bra{\perp}\ket{\perp} \\
    \tilde{\varsigma} &= \bra{1}\ket{x_{\kbar}} \\
                      &= \frac{\alpha}{\sqrt{\bra{1}\ket{1}}} \bra{1}\ket{1} + \bra{1}\ket{\perp} \\
                      &= \alpha\sqrt{\bra{1}\ket{1}} \\
                      &= \alpha\sqrt{d-1}
\end{align*}
The integral then becomes:
\begin{align*}
    I &= \int_{\mathbb{R}}\int_{\mathbb{R}^{d-2}} \exp\qty(-a\qty[\alpha^2 + \bra{\perp}\ket{\perp} - b\alpha^2\qty(d-1) - 2\alpha\sqrt{d-1}\qty(t + s - 2\mu_k)]) \dd{\alpha}\dd[n-2]{\ket{\perp}} \\
        &= \int_{\mathbb{R}} \exp\qty(-a\qty[\alpha^2 - b\alpha^2\qty(d-1) - 2\alpha\sqrt{d-1}\qty(t + s - 2\mu_k)]) \dd{\alpha} \times \\
        &\quad \times \int_{\mathbb{R}^{d-2}} \exp(-a\bra{\perp}\ket{\perp}) \dd[n-2]{\ket{\perp}} \\
        &= \qty(\frac{\pi}{a})^{\frac{d-2}{2}} \exp(\frac{a\qty(d-1)\qty(t + s - 2\mu_k)^2}{\qty(1 - b\qty(d-1))})\int_{\mathbb{R}} \exp(-a\qty(\alpha - \frac{2a\sqrt{d-1}}{2a\qty(1 - b\qty(d-1))})^2) \dd{\alpha} \\
        &= \qty(\frac{\pi}{a})^{\frac{d-1}{2}} \exp(\frac{a\qty(d-1)\qty(t + s - 2\mu_k)^2}{\qty(1 - b\qty(d-1))}) \\
        &= \qty(\pi\qty(1-r))^{\frac{d-1}{2}} \exp(\frac{\qty(1 + r\qty(d-1))\qty(d-1)\qty(t + s - 2\mu_k)^2}{1-r})
\end{align*}
Putting everything together, we get:
\begin{align*}
    \kappa_{\kbar}\qty(t, s) &= \frac{\exp(-\frac{1 + r\qty(d - 2)}{2\qty(1-r)\qty[1 + r\qty(d-1)]}\qty[\qty(t - \mu_k)^2 + \qty(s - \mu_k)^2])}{\qty(2\pi)^{d}\qty(1 + r\qty(d-1))\qty(1-r)^{d-1}} \times \\
                             &\quad \times \qty(\pi\qty(1-r))^{\frac{d-1}{2}} \exp(\frac{\qty(1 + r\qty(d-1))\qty(d-1)\qty(t + s - 2\mu_k)^2}{1-r}) \\
                             &= \frac{\qty(1-r)^{\frac{1-d}{2}}}{2^{d} \pi^{\frac{1+d}{2}}\qty(1 + r\qty(d-1))} \exp(-\frac{\qty(1 + r\qty(d - 2))\qty[\qty(t - \mu_k)^2 + \qty(s - \mu_k)^2]}{2\qty(1-r)\qty(1 + r\qty(d-1))}) \times \\
                             &\quad \times \exp(\frac{\qty(1 + r\qty(d-1))\qty(d-1)\qty(t + s - 2\mu_k)^2}{1-r})
\end{align*}
We can write this simply as:
\begin{align*}
    \kappa_{\kbar}\qty(t, s) &= \mathcal{C} \exp\qty(-u\qty(t^2 + s^2) + vts)
\end{align*}
for some real constants \(\mathcal{C}\), \(u\) and \(v\). The eigenfunctions of
an integral operator with a kernel of this form are known to be\footnote{This is
obviously a fact known to everyone, and not something a machine that does a
bunch of matrix vector products that is very good at text prediction wrote.} the
Hermite functions. The idea will be to use the Hermite polynomial's
orthogonality,~\cref{eq:ortho}, and Mehler's formula~\cref{eq:mehler}
(\href{https://en.wikipedia.org/wiki/Mehler_kernel?oldformat=true#Physics_version}{found
on wikipedia}) to recover an eigenvalue equation.
\begin{align*}
    \int_{-\infty}^{\infty} H_n\qty(x)H_m\qty(x) e^{-x^2} \dd{x} &= 2^n n! \sqrt{\pi} \delta_{nm}\stepcounter{equation}\tag{\theequation}\label{eq:ortho} \\
    \frac{1}{\sqrt{1 - \rho^2}}\exp(\frac{-\qty(1 + \rho^2)\qty(x^2 + y^2) + 4xy\rho}{2\qty(1 - \rho^2)}) &= \sum_{n=0}^{\infty} \qty(\frac{\rho}{2})^n \frac{H_n\qty(x)H_n\qty(y)}{n!} \exp(-\frac{x^2 + y^2}{2})\stepcounter{equation}\tag{\theequation}\label{eq:mehler}
\end{align*}

\todo{just use hermite functions and you're good}

Recall that we are after the eigenfunctions \(v_k^{\qty(n)}\qty(s)\) (the right
eigenvectors of \(F_{\qty(k)}\)) corresponding eigenvalues
\(\lambda_k^{\qty(n)}\)) such that \(\displaystyle\qty(C_kv_k^{\qty(n)})\qty(s)
= \int_{\mathbb{R}} \kappa_{\kbar}\qty(t, s) v_k^{\qty(n)}\qty(t) \dd{t} =
\lambda_k^{\qty(n)}v_k^{\qty(n)}\). Starting from~\cref{eq:mehler} and working
towards~\cref{eq:ortho}, we can get:
\begin{align*}
    \int_{\mathbb{R}} \frac{1}{\sqrt{1 - \rho^2}} \exp(\frac{-\qty(\frac{3}{2} + \rho^2)\qty(x^2 + y^2) + 4xy\rho}{2\qty(1 - \rho^2)}) H_n\qty(x) \dd{x} &= \sum_{n=0}^{\infty} \qty(\frac{\rho}{2})^n \frac{H_n\qty(y)e^{-y^2}}{n!} \int_{\mathbb{R}}  H_n\qty(x)H_n\qty(x) e^{-x^2} \dd{x} \\
                                                                                                                                                         &= \sum_{n=0}^{\infty} \qty(\frac{\rho}{2})^n H_n\qty(y) e^{-y^2}
\end{align*}
% If we can identify the kernel \(\kappa_{\kbar}\qty(t, s)\) with the kernel in
% the left-hand side of the above equation, then we can read off the
% eigenfunctions and eigenvalues of \(C_k\) from the right-hand side. We can do
% this by setting \(\rho = 2\sqrt{uv}/(u + v)\) and rescaling \(t\) and \(s\) by
% \(\sqrt{(u + v)/2}\). This gives us:

\end{document}
